% ============================================================
%  VII. RESULTS AND DISCUSSION
% ============================================================
\section{Results and Discussion}
\label{sec:results}

\subsection{Search Time Comparison}

Table~\ref{tab:search_time} presents the search time results from the 50
physical trials. Across all five patient yaw angles, the Spot+Arm
configuration achieves a mean search time of \SI{9.4}{\second}, compared to
\SI{15.5}{\second} for the Orbbec-Only configuration, yielding a
\textbf{1.6}$\times$ mean speedup. The median speedup is \textbf{2.0}$\times$,
indicating that for the majority of patient orientations the arm provides a
substantial advantage.

\begin{table}[htbp]
  \centering
  \caption{Search time comparison: Spot+Z1 arm vs.\ Orbbec-only camera,
           measured over 50 physical trials (five per angle per
           configuration).  Values are mean $\pm$ standard deviation.}
  \label{tab:search_time}
  \begin{tabular}{lcc}
    \toprule
    & \textbf{Spot+Arm} & \textbf{Orbbec-Only} \\
    \textbf{Patient Yaw} & \textbf{Time (s)} & \textbf{Time (s)} \\
    \midrule
    $0\degree$   & $1.4 \pm 0.4$  & $3.2 \pm 0.6$ \\
    $45\degree$  & $5.2 \pm 1.1$  & $13.4 \pm 2.0$ \\
    $90\degree$  & $8.0 \pm 1.6$  & $13.8 \pm 2.4$ \\
    $135\degree$ & $10.8 \pm 2.2$ & $23.5 \pm 3.8$ \\
    $180\degree$ & $21.8 \pm 3.5$ & $23.8 \pm 4.0$ \\
    \midrule
    \multicolumn{3}{c}{\textbf{Aggregate}} \\
    \midrule
    Mean & 9.4 & 15.5 \\
    \midrule
    \textbf{Speedup (O/S)} & \multicolumn{2}{c}{1.6$\times$ (mean), 2.0$\times$ (median)} \\
    \bottomrule
  \end{tabular}
\end{table}

At individual yaw angles, the speedup ranges from \textbf{2.6}$\times$ at
$45\degree$ (\SI{5.2}{\second} vs.\ \SI{13.4}{\second}) to \textbf{1.1}$\times$
at $180\degree$, where both configurations require a full base rotation to
the opposite side. At $0\degree$ (patient facing the robot), the arm provides
a \textbf{2.3}$\times$ speedup (\SI{1.4}{\second} vs.\ \SI{3.2}{\second}),
as the \ac{EE}-mounted camera can immediately detect the patient without
waiting for the base to rotate. At $90\degree$ and $135\degree$, the speedup
is \textbf{1.7}$\times$ and \textbf{2.2}$\times$ respectively, demonstrating
consistent advantage across the frontal hemisphere. The arm's benefit is
strongest between $0\degree$ and $135\degree$ and diminishes at $180\degree$,
where the patient lies opposite the robot's initial heading and both
configurations must sweep through nearly the full azimuth.

In absolute terms, the Orbbec-Only configuration shows larger standard
deviations at mid-range angles (\SI{2.0}{\second}--\SI{4.0}{\second} vs.\
\SI{1.1}{\second}--\SI{3.5}{\second} for Spot+Arm), reflecting the
cumulative timing uncertainty from a larger number of base rotations.
However, its coefficients of variation are lower
(\SI{15}{\percent}--\SI{19}{\percent}) than those of Spot+Arm
(\SI{16}{\percent}--\SI{29}{\percent}) because the Orbbec-Only search time
grows proportionally with the number of positions visited, keeping the
relative spread narrow.  The Spot+Arm configuration, by contrast, exhibits
higher relative variance because detection within the arm's 7-pose scan
depends on the angular offset from the scan centre---a patient near the edge
of the coverage window is found late in the scan, introducing proportional
timing uncertainty that can nearly double the search time relative to a
patient near the centre.

Overall, the \textbf{1.6}$\times$ mean speedup means that \ac{TERESA}
reduces the expected time to localise a patient at an unknown orientation by
approximately \SI{40}{\percent}, from \SI{15.5}{\second} to
\SI{9.4}{\second}. This reduction is achieved without any additional sensing
hardware: the Z1 arm, already required for ultrasound probe manipulation,
serves a dual role as an active perception platform during the search phase.
In an emergency response scenario, the \SI{6}{\second} saving is critical,
as it directly shortens the interval between robot arrival and the start of
the medical assessment.

\subsection{Additional Results}

\TODO{Body scanning accuracy and completeness results}

\TODO{FAST ultrasound probe positioning accuracy}

\TODO{Ablation study of perception-quality-aware velocity control}
